
%%%%%%%%%%%%%%%%%%%%%%%%%%%%%%%%%%%%%%%%%
%     ACA 2013 Malaga - Edu Session     %        
%     Yo, Justo                         %
%%%%%%%%%%%%%%%%%%%%%%%%%%%%%%%%%%%%%%%%%

\documentclass[a4paper]{article} %
\usepackage{graphicx,amssymb} %

\usepackage{url}
\usepackage{color}

\textwidth=15cm \hoffset=-1.2cm %
\textheight=25cm \voffset=-2cm %

\pagestyle{empty} %

\date{} %

\def\keywords#1{\begin{center}{\bf Keywords}\\{#1}\end{center}} %

\def\titulo#1{\title{#1}} %
\def\autores#1{\author{#1}} %

% Please, do not change any of the above lines

\begin{document}

% Type down your paper title
\titulo{Some maths problems for the average citizen}

% Authors
\autores{Eugenio Roanes-Lozano \\ %  Affiliation
       Instituto de Matem\'atica Interdisciplinar (IMI),\\
       Departamento de \'Algebra, Facultad de Educaci\'on, \\
       Universidad Complutense de Madrid, E-28040 Madrid (Spain) \\ \\
       Justo Cabezas-Corchero\\
       High School Mathematics Professor (Ret.)\\
       Badajoz, Spain\\ \\
       \tt{eroanes@mat.ucm.es} % Only one corresponding e-mail
       }%


\maketitle

\thispagestyle{empty}



% The abstract

\begin{abstract}
We have been always surprised by the inability of many students
to adequately treat some problems of everyday life that only
require of elementary mathematics (meanwhile they can solve
problems with a much more complicated background). Obviously, the
reason is that they haven't been trained in facing these sorts of
problems. We shall give an overview of some of the
typical mathematical problems that an average citizen needs to
know how to face. In most of them a CAS helps (avoiding the
tedious computations), but cannot solve the problem by itself.
But we believe that CAS can be key in the process of
experimenting with these problems in order to achieve the
necessary skills for solving them in real life.
\end{abstract}

\keywords{Basic Mathematics, Curriculum, CAS} % Write down at least 3 Keywords


\section{Introduction}

The first author has taught at the School of Statistics, the
School of Education and the School of Mathematics of the
Universidad Complutense de Madrid for more than 25 years. His
students have ranged from freshmen to Ph.D. students.

The second author has been a mathematics high school teacher
for more than 35 years. He ha also taught maths and statistics
at different schools of the Universidad de Extremadura.

We have been always surprised by the inability of many students to
adequately treat some problems of everyday life that only require
of elementary mathematics (meanwhile they can solve problems with
a much more complicated background). Obviously the reason is that
they haven't been trained in facing these sorts of problems.

An example are percentages. The first author is now teaching a
subject entitled ``Elementary mathematics with computer'' to 2nd
year pedagogy degree students (in it, elementary mathematical
problems are treated using CAS and DGS). While most these
students could calculate a discount on a price, not all could
correctly solve a similar problem: write a computer program that,
given the total diameters of two tires, decides whether the new
tire differs in more than a 3 \% w.r.t. the old one (the obstacle
was mathematical, not computer related). Moreover,
they were not sure if adding a tax plus a discount to a price
(both expressed as percentages) did commute. Meanwhile, they use
SPSS in complicated statistical studies regarding pedagogical
issues.

In ancient cultures mathematical research was mainly focused on
topics with a direct application: plane geometry (for surveying
and architecture) and astronomy (for religious reasons), although
this knowledge was not intended for the average citizen. But
these interests have been changing in the last thousands years.

Many of the problems that could now be of ``general
interest'' have an economical background, and many require the
use of elementary physics (mainly physical units). Some examples
can be found afterwards.



\section{Some of the examples detected}

Examples of some of the topics that we have detected
that could be labeled as \textit{important} for an average
citizen (\textit{essential skills}) can be found afterwards. We
have only included situations that we believe are not correctly
covered (in practice) by the Spanish educational system.

\medskip
\noindent\underline{Economy:}

\begin{itemize}
\item Do you save money if you change your present (perfectly working)
      refrigerator by an A++ one?
\item According to your present telephone invoice, should you migrate to
      another company with a completely different way to bill the phone
      calls that looks much cheaper?
\item Should you change your present car (in its midlife) by a much
      more ecological new hybrid for economical reasons?
\item Do low consumption bulbs worth their higher price?
\end{itemize}

Nevertheless, there are many other mathematical topics related to real life
situations. Some examples are include afterwards.

\medskip
\noindent\underline{Divisibility:}

\begin{itemize}
\item My living room is $5.40m \times 4.20m$ and I would like to floor it with
      tiles but I don't have a cut tile machine. Which is the maximum tile size
      that I can use?
\end{itemize}

\medskip
\noindent\underline{Dilating areas and volumes and scales:}
\begin{itemize}
\item This frustoconical drinking glass looks a bit too small. I'll buy this
      other one that is $1.2$ times higher and wider. It is just a bit bigger
      than the other one, right?
\item The rooms in this house floor plan look huge.
      Is the furniture represented in the house floor plan in the correct scale?
\end{itemize}

\medskip
\noindent\underline{Derivatives (or increments):}
\begin{itemize}
\item I can find in a local newspaper that ``The unemployment has lowered
      its growth'' meanwhile another paper (of a different political tendency)
      publishes the same day that ``The unemployment has grown''.
      Can we assure that, at least, one of them is lying?
\end{itemize}

\medskip
\noindent\underline{Combinatorics:}
\begin{itemize}
\item In a certain lottery there are 100 numbers. The probability to win if
      I buy one day one number is 1\%. If I buy a ticket today and a ticket
      tomorrow the probability to win is 1\%+1\%=2\%. OK?
\end{itemize}

\medskip
\noindent\underline{Cardinal of the union / Probability of the union:}
\begin{itemize}
\item In this class there is a $45\%$ of males and a $30\%$ of people from
      Andalusia. Do you agree that one way or another we cover $75\%$ of
      the students?
\end{itemize}

\medskip
\noindent\underline{Negative exponential:}
\begin{itemize}
\item Why should I take antibiotics (and many other drugs) following a strict
      dosage schedule?
\item I know that Carbon-14 is an unstable isotope of Carbon, but how Carbon-14
      dating method works?
\end{itemize}


\medskip
\noindent\underline{Correlation:}
\begin{itemize}
\item There is no ``functional relation'' between the height and weight of people,
      but, is there any other kind of ``relation''?
\end{itemize}

\medskip
\noindent\underline{Normal distribution:}
\begin{itemize}
\item How do shops decide which shoe sizes should they have in stock?
\item For which range of IQ students is curriculum designed?
\end{itemize}

\medskip
\noindent\underline{Interpretation of graphic representations:}
\begin{itemize}
\item In the figure below %Figure~1
      the evolution a certain stock exchange index along one
      week is represented. It reflects huge economical movements, doesn't it?
\end{itemize}

%***************** FIGURA **************
\begin{figure}[htb]
\begin{center}
\scalebox{.5}{
\includegraphics  %[width=\textwidth]
{Fig1.jpg}
} %
%\caption{Evolution a certain stock exchange index.}
\end{center}
\end{figure}
%***************************************

Some mathematical facts have even been incorporated into the
acquis communautaire:

\medskip
\noindent\underline{Partial ordering:}
\begin{itemize}
\item According to a Spanish saying: ``All comparisons are obnoxious''
      (``Todas las comparaciones son odiosas'').
\end{itemize}



\section{Remark}

Once the first draft of this Extended Abstract was already
prepared, one of the main Spanish newspapers published an
impressive report about the failures in elementary mathematics
(and other subjects) of students with a degree in Primary School
Teaching during their competitive recruitment examinations at
Madrid region\footnote{Pilar \'Alvarez, Maestros suspensos en
primaria (in Spanish).
                          El Pa\'{\i}s, March 13th 2013
                          \url{http://sociedad.elpais.com/sociedad/2013/03/13/actualidad/1363202478_209351.html}.\\
                          Anonymous, Las matem\'aticas se resisten (in Spanish).
                          El Pa\'{\i}s, March 13th 2013.
                          \url{http://sociedad.elpais.com/sociedad/2013/03/13/actualidad/1363201114_422663.html}
                         }.
The exercise with the worse results (7.09\% of the answers were
correct) was an elementary example about time, weight and area
unit conversion.



\section{Conclusions}

Taking into account that most citizens are not mathematicians, we
believe that this sort of topics/problems, closer to
everyday life, should be included into the curricula, in order
to provide a ``more useful'' mathematics education.

Moreover, the computations required by many of these problems
can be bypassed using a CAS (CAS can even carry units along the computations).

An issue to be discussed is whether it would be
better to treat these topics/problems within the traditional
curricula or in separated workshops.



\end{document}
